\section{Compact random ideal sampling algorithms}
\label{sec:sampling}

After replacing HNF algorithm by \textsf{MLLL} in the ideal operations,
the ideal operations no longer dominate the fixed-precision budget.
Instead, the largest remaining intermediate values arise from the prime-norm branch of \textsf{RandomIdealGivenNorm}.
Also, the previous analysis of the algorithm \textsf{RandomEquivalentPrimeIdeal} uses a conservative bound on the norm of the sampled equivalent ideal.

We therefore revisit these routines.
In subsection~\ref{subsec:RandomIdealGivenNorm}, we introduce the modified algorithm of \textsf{RandomIdealGivenNorm}, by reducing the maximum size of intermediate values of the algorithm.
In subsection~\ref{subsec:RandomEquivalentPrimeIdeal}, we give a tighter size analysis.
These refinements improve the final
uniform fixed-precision bounds during the key generation and signing procedures of SQIsign.

% In this section, we modify/analyze two other (ideal sampling) algorithms to reduce the size of integers in the execution of KeyGen/Sign of SQIsign.

% In this section, we revisit two ideal sampling algorithms used in SQIsign KeyGen and Sign procedures.
% We preserve the output distribution while reducing the size of the integers manipulated during the computation.

\subsection{Modification of RandomIdealGivenNorm}
\label{subsec:RandomIdealGivenNorm}

We first consider the prime-norm branch of \textsf{RandomIdealGivenNorm} algorithm~\cite{AAA+25}. In the original algorithm, after sampling $\gamma$ and $\beta$, one returns
\[
J'=\mathcal O_0(\gamma\beta)+N\mathcal O_0.
\]
However, this ideal depends only on the residue class of $\gamma\beta$ modulo $N\mathcal O_0$. There is therefore no need to materialize the unreduced product: it suffices to compute a representative
\[
g \equiv \gamma\beta \pmod{N\mathcal O_0}
\]
and return $\mathcal O_0 g + N\mathcal O_0$. In particular, the final multiplication may be carried out modulo $N$, so the coefficients handled by the algorithm remain reduced modulo $N$.

The same observation also simplifies the sampling of $\beta$.
In the original algorithm, the rejection condition is exactly $\nrd(\beta)\not\equiv 0 \pmod N$.
Since only $\beta \bmod N\mathcal O_0$ is used afterwards, we may sample $\beta$ directly from residue classes with nonzero reduced norm modulo $N$.
This yields the compact variant introduced in Algorithm~\ref{alg:RandomIdealGivenNorm}.

\begin{algorithm}
\caption{RandomIdealGivenPrimeNorm($N$)}
\label{alg:RandomIdealGivenNorm}
\begin{algorithmic}[1]
\Statex \textbf{Input:} A prime number $N > 3$.
% satisfying $N\equiv1\pmod4$
% which is the norm of some left $\mathcal{O}_0$ ideal.
\Statex \textbf{Output:} A random left $\calO_0$-ideal $J'$ of norm $N$.

% \State $(x,y)\gets\mathbf{Cornacchia}(1,N)$
% \State $(z,t)\gets\mathbf{Cornacchia}(1,N)$
\State $\gamma \gets 0+0i+0j+0k$
\While{$\gcd\left(\nrd(\gamma), N^2\right) = N$}
    \State Sample integers $g_1,g_2,g_3 \sample [0, N-1]$ uniformly
    \State $\gamma \gets g_1 i + g_2 j + g_3 ij$
    \If{$(-\nrd(\gamma):N) = 1$}
        \State $\gamma \gets \gamma + \left(\sqrt{-\mathrm{nrd}(\gamma)} \bmod N\right)$
        \EndIf
    \EndWhile
\State Sample integers $a,b,c,d \sample [0,N-1]$ s.t. $a^2+b^2+p(c^2+d^2)\neq 0 \pmod N$ uniformly
\State $\beta \gets a+bi+cj+dk$
\State $g \gets \gamma\beta \bmod N$
\State $J' \gets$ the left $O_0$-ideal generated by $g$ and $N$
\State \Return $J'$
\end{algorithmic}
\end{algorithm}

The following lemma isolates the fact needed for correctness: reducing a generator modulo $N\mathcal O_0$ preserves norm divisibility by $N$, so it does not change the generated ideal.

\begin{lemma}
\label{lem:nrd-mod}
Let $N$ be a prime, and let $\beta,\gamma \in \mathcal O_0$ satisfy
\[
\gcd\!\bigl(N^2,\nrd(\gamma)\bigr)=N,
\qquad
\gcd\!\bigl(\nrd(\beta),N\bigr)=1.
\]
Set $\alpha := \gamma\beta$.
Let $\widetilde{\alpha}\in \mathcal O_0$ be any lift of the residue class
$\alpha \bmod N\mathcal O_0$, i.e.
\[
\widetilde{\alpha}\equiv \alpha \pmod{N\mathcal O_0}.
\]
Then:
\begin{enumerate}
    \item
    $
    \gcd\!\bigl(N^2,\nrd(\alpha)\bigr)=N;
    $
    \item
    $
    \nrd(\widetilde{\alpha})
    \equiv
    \nrd(\alpha)
    \pmod N.
    $
    In particular,
    $N \mid \nrd(\widetilde{\alpha})$;
    \item
    $\mathcal O_0\langle \widetilde{\alpha},N\rangle
    =
    \mathcal O_0\langle \alpha,N\rangle.$
\end{enumerate}
\end{lemma}

\begin{proof}
See Appendix~\ref{subsec:nrd-mod}.
\end{proof}

% \begin{corollary}
% With the same notation, if
% \[
% I=\mathcal O_0\langle \alpha,N\rangle
% =
% \mathcal O_0\langle \widetilde{\alpha},N\rangle,
% \]
% then on $E_0[N]$ one has $\widetilde{\alpha}=\alpha$, and therefore
% \[
% E_0[I]
% =
% \ker(\alpha)\cap E_0[N]
% =
% \ker(\widetilde{\alpha})\cap E_0[N].
% \]
% \end{corollary}

% \begin{proof}
% For every $P\in E_0[N]$,
% \[
% \widetilde{\alpha}(P)
% =
% (\alpha+N\delta)(P)
% =
% \alpha(P)+\delta(NP)
% =
% \alpha(P),
% \]
% since $NP=0$.
% \end{proof}

\paragraph{Uniform distribution of Algorithm~\ref{alg:RandomIdealGivenNorm}.}
The sampling of $\gamma$ is unchanged from the original prime-norm branch of
\textsf{RandomIdealGivenNorm}. For $\beta$, the original algorithm samples
coordinates in $[1,N]$ and rejects unless $\gcd(\nrd(\beta),N)=1$, whereas
Algorithm~\ref{alg:RandomIdealGivenNorm} samples directly from
$[0,N-1]^4$ subject to
\[
\nrd(\beta)\not\equiv 0 \pmod N.
\]
Since $N$ is prime, these two conditions are equivalent
% , and coordinate-wise
% reduction modulo $N$ gives a bijection between $[1,N]^4$ and $[0,N-1]^4$
. Hence both algorithms induce the same distribution on the residue class of
$\beta$ modulo $N\mathcal O_0$.

It therefore remains only to justify the modification in the last step:
the original routine constructs the ideal from the unreduced product
$\gamma\beta$, whereas Algorithm~\ref{alg:RandomIdealGivenNorm} uses a
representative of $\gamma\beta \bmod N\mathcal O_0$. The next lemma proves
the stronger pointwise statement that, for every realization of the random
choices, both constructions return the same left $\mathcal O_0$-ideal.
Equality of output distributions follows immediately.

\begin{lemma}
\label{lem:RandomIdealGivenNorm-distribution}
Let $N$ be a prime and $\beta,\gamma \in \calO_0$ such that $N \mid \nrd(\gamma)$, $N \nmid \nrd(\beta)$.
Define
\[
J'_{\mathrm{old}}:=\mathcal O_0(\gamma\beta)+N\mathcal O_0,
\qquad
J'_{\mathrm{new}}:=\mathcal O_0(\overline{\gamma\beta})+N\mathcal O_0,
\]
where $\overline{\gamma\beta}$ is any lift in $\mathcal O_0$ of the class of
$\gamma\beta$ modulo $N\mathcal O_0$.
% (equivalently, in an implementation, the reduction of the coordinates of $\gamma\beta$ modulo $N$ in the fixed integral basis).
Then
\[
J'_{\mathrm{old}}=J'_{\mathrm{new}}.
\]
In particular, replacing $\gamma\beta \bmod N$ by $\gamma\beta$ in the last line of
Algorithm~\ref{alg:RandomIdealGivenNorm} does not change the output distribution.
\end{lemma}

\begin{proof}
See Appendix~\ref{subsec:RandomIdealGivenNorm-distribution}.
% More generally, if $a,a'\in \mathcal O_0$ satisfy $a-a'\in N\mathcal O_0$, then
% \[
% \mathcal O_0 a+N\mathcal O_0=\mathcal O_0 a'+N\mathcal O_0.
% \]
% Indeed, writing $a=a'+N\delta$ with $\delta\in\mathcal O_0$, we get
% \[
% \mathcal O_0 a \subseteq \mathcal O_0 a' + N\mathcal O_0,
% \]
% hence
% \[
% \mathcal O_0 a + N\mathcal O_0 \subseteq \mathcal O_0 a' + N\mathcal O_0.
% \]
% The reverse inclusion follows symmetrically from $a'=a-N\delta$.

% Now take $a=\gamma\beta$ and $a'=\overline{\gamma\beta}$. By construction,
% \[
% \gamma\beta-\overline{\gamma\beta}\in N\mathcal O_0.
% \]
% Therefore
% \[
% \mathcal O_0(\gamma\beta)+N\mathcal O_0
% =
% \mathcal O_0(\overline{\gamma\beta})+N\mathcal O_0.
% \]
% Thus, for every realization of the random choices, the modified and the original
% algorithms return exactly the same left $\mathcal O_0$-ideal. Consequently their
% output distributions are identical.
\end{proof}

% \begin{remark}
% We heuristically assert that the distribution of outputs of Algorithm~\ref{alg:RandomIdealGivenNorm} is statistically indistinguishable with uniform distribution, by experiment. See Appendix~\ref{subsec:exp-RandomIdealGivenNorm}.
% \end{remark}

% \begin{algorithm}[H]
% \caption{CompactRandomIdealGivenNorm($N$,prime)}
% \label{alg:RandomIdealGivenNorm}
% \begin{algorithmic}[1]
% \Statex \textbf{Input:} A positive integer $N\ge2$ not multiple of $p$ which is the norm of some left $\mathcal{O}_0$ ideal, A boolean \emph{prime} indicating whether $N$ is prime.
% \Statex \textbf{Output:} A random left ideal $J'$ of $\calO_0$ of norm $N$, or raise an exception.

% \State $found \gets \text{false}$
% \If{prime}
%   \While{not $found$}
%     \State $g_1,g_2,g_3 \gets$ independent uniformly random integers in $[0, N-1]$
%     \State $\gamma \gets g_1 i + g_2 j + g_3 k$
%     \State $found \gets \bigl(1=\mathrm{Legendre}(-\,\mathrm{nrd}(\gamma),N)\bigr)$
%     \If{$found$}
%       \State $\gamma \gets \gamma + \mathrm{ModularSQRT}(-\,\mathrm{nrd}(\gamma),N)$
%     \EndIf
%   \EndWhile
% \Else
%   \State $k\gets\textbf{NextPrime}(p+1)-p$
%   \State $g_1,g_2\gets\mathbf{Cornacchia}(1,kN)$
%   \State $g_3,g_4\gets\mathbf{Cornacchia}(1,N)$
%   \State $\gamma\gets g_1+g_2i+g_3j+g_4k$
%   % \If{$N>p$}
%   %   \State $\gamma \gets$ \textbf{GeneralizedRepresentInteger}$(N, i, O_0, \text{false})$
%   % \Else
%   %   \State $m \gets \lceil\frac{p}{N}\rceil$
%   %   \State $\gamma \gets$ \textbf{GeneralizedRepresentInteger}$(mN, i, O_0, \text{false})$
%   % \EndIf
% \EndIf
% \While{not $found$}
%   \State $x,y,z,w \gets$ uniformly randomly selected integers in $[1,N]$
%   \State $\beta \gets x + y i + z j + w ij$
%   \State $found \gets \bigl(\gcd(\mathrm{nrd}(\beta),N)=1\bigr)$
% \EndWhile
% \State $J' \gets$ the left $O_0$-ideal generated by $\gamma\beta$ and $N$
% \State \Return $J'$
% \end{algorithmic}
% \end{algorithm}

\begin{lemma}
\label{lem:RandomIdealGivenNorm-bound}
During the execution of Algorithm~\ref{alg:RandomIdealGivenNorm}, the maximum size of integers is less than $4N^2$.
\end{lemma}

\begin{proof}
\noindent\textbf{(1) Lines 1--3.}
During the execution of \textsf{Cornacchia}, the maximum size of integers does not grow by~\cite[Lemma~18]{KLKL25}.
Hence, coefficients representing $\gamma$ are less than $\sqrt{N}$.

\smallskip
\noindent\textbf{(2) Lines 4--6.}
Since $a,b,c,d < N$, $\beta$ has coefficients less than $N$.
Hence, in line~6, $\gamma\beta$ has coefficients less than $N$.
Computing the ideal $\calO_0(\gamma\beta)+\calO_0 N$ is as follows:
\begin{enumerate}
    \item multiply $\calO_0=\langle 1,i,\frac{1+k}{2},\frac{i+j}{2} \rangle$ by $\gamma\beta$ and encode it to a matrix with \[\left(1\cdot(\gamma\beta), i\cdot(\gamma\beta), \frac{1+k}{2}\cdot(\gamma\beta), \frac{i+j}{2}\cdot(\gamma\beta)\right).\] The maximum size of integers to represent each coefficient is less than $N\sqrt{N}$.
    \item By the same procedure, multiply $\calO_0$ by $N$. The maximum size of integers to represent each coefficient is less than $N$.
    \item Concatenate two vectors and then apply Algorithm~\ref{alg:ML2}, ensuring the output matrix is a basis matrix of $I=\calO_0(\gamma\beta)+\calO_0 N$. In this procedure, the maximum size of integers is less than $4N^3$ by Lemma~\ref{lem:ML2-bound}.
\end{enumerate}
\end{proof}


\subsection{Output norm of RandomEquivalentPrimeIdeal}
\label{subsec:RandomEquivalentPrimeIdeal}

We recall the crude worst-case bound for \textsf{RandomEquivalentPrimeIdeal} algorithm used in previous analyses~\cite{KLKL25}.
% Let
% \[
%   \beta=\sum_{i=1}^{4} c_i\alpha_i,\quad |c_i|\le \coeffbound,
% \]
% where $\alpha_1,\ldots,\alpha_4$ is the LLL-reduced basis of the ideal lattice $I$.
% Since the reduced norm is a positive definite quadratic form, the triangle inequality gives
% \[
%   \sqrt{\nrd(\beta)}
%   \leq \sum_{i=1}^{4} |c_i|\sqrt{\nrd(\alpha_i)}
%   \leq \coeffbound\sum_{i=1}^{4}\sqrt{\nrd(\alpha_i)}.
% \]
% Thus, for $J=\chi_I(\beta)$, one obtains the rough bound
% \[
%   \nrd(J)
%   = \frac{\nrd(\beta)}{\nrd(I)}
%   \leq
%   \frac{\coeffbound^2}{\nrd(I)}
%   \left(\sum_{i=1}^{4}\sqrt{\nrd(\alpha_i)}\right)^2 .
% \]
The bound on output norm was deriven in a coarse empirical manner:
one directly observed the norm of the ideal returned by Algorithm~\ref{alg:RandomEquivalentPrimeIdeal} and chose a sufficiently large global upper bound as $p$, while the observed output norm looks $O(\sqrt{p})$.
In contrast, we heuristically compute a tight bound on the size of $\nrd(\alpha_4)$.
This allows the final bound on $\nrd(J)$ to remain empirical in nature, but to be much more tightly calibrated to the actual sampling process.

\begin{algorithm}
\caption{\textsf{RandomEquivalentPrimeIdeal}($I$)}
\label{alg:RandomEquivalentPrimeIdeal}
\textbf{Input:} $I$, a left $\mathcal{O}$-ideal.\\
\textbf{Output:} $J \sim I$ of small prime norm, or raise an exception if unsuccessful.
\begin{algorithmic}[1]
\State Initialize $\textit{counter} \gets 0$
\State $(\alpha_{1},\alpha_{2},\alpha_{3},\alpha_{4}) \gets \mathsf{LLL}(I;\delta=\frac34)$
\While{$\textit{counter} < (2\cdot \texttt{QUAT\_equiv\_bound\_coeff} + 1)^{4}$}
  \State $\textit{counter} \gets \textit{counter} + 1$
  \State Sample $c_{1},c_{2},c_{3},c_{4}$ uniformly at random from $[-b,\ldots,b]$, for bound $b=\texttt{QUAT\_equiv\_bound\_coeff}$
  \State $\beta \gets \sum_{i=1}^{4} c_{i}\alpha_{i}$
  \State $J \gets \chi_{I}(\beta)$
  \If{$\nrd(J)$ is prime}
     \State \Return $J$
  \EndIf
\EndWhile
\State \textbf{raise} \textsc{Exception}(``RandomEquivalentPrimeIdeal failed'')
\end{algorithmic}
\end{algorithm}

\begin{lemma}
\label{lem:RandomEquivalentPrimeIdeal-bound}
The reduced norm of output ideal of Algorithm~\ref{alg:RandomEquivalentPrimeIdeal} is less than $2^{33.5}\sqrt{p}$ when given an input ideal generated by Algorithm~\ref{alg:RandomIdealGivenNorm} with given norm $\Dmix$.
\end{lemma}

\begin{proof}
Let $I$ be the output of Algorithm~\ref{alg:RandomIdealGivenNorm} and $(\alpha_1,\alpha_2,\alpha_3,\alpha_4)$ be its LLL-reduced basis as in line~2.
By combining Lemma~\ref{lem:transference}, Heuristic~\ref{heu:dual}, and Lemma~\ref{lem:lll-successive-minima}, we have
\[ \frac{\nrd(\alpha_4)}{\nrd(I)} \le 2^{13.5}\cdot\sqrt{p}. \]
Hence, we have
\begin{align*}
    \nrd(J)
    &\le \frac{\nrd\left(\coeffbound\cdot\sum\limits_{i=1}^{4} \alpha_i\right) }{\nrd(I)}\\
    &\le \frac{\coeffbound^2\cdot\nrd(4\alpha_4)}{\nrd(I)}\\
    &= \frac{16\cdot\coeffbound^2\cdot\nrd(\alpha_4)}{\nrd(I)}\\
    &\le 2^{16}\cdot2^{13.5}\cdot\sqrt{p} = 2^{29.5}\cdot\sqrt{p}
    \enspace .
\end{align*}
\end{proof}

Heuristic~\ref{heu:dual} is the heuristic which states the lower bound of successive minima of the trace-dual lattice of the ideal sampled by Algorithm~\ref{alg:RandomIdealGivenNorm}.
By using the transference theorem, we get the tight bound of the output norm of Algorithm~\ref{alg:RandomEquivalentPrimeIdeal}.

\begin{heuristic}
\label{heu:dual}
Let $I$ be the output of Algorithm~\ref{alg:RandomIdealGivenNorm} with given input norm $N$ and $(\alpha^\#_1,\alpha^\#_2,\alpha^\#_3,\alpha^\#_4)$ be the Minkowski-reduced basis of its trace-dual lattice.
Then, \[N\cdot\nrd(\alpha^\#_1) \ge \frac{2^{-8}}{\sqrt{p}}\enspace.\]
\end{heuristic}

\begin{proof}
See Appendix~\ref{subsec:exp-dual}.
\end{proof}

% Let $I$ be a left $\mathcal O_0$-ideal and let
% \[
% N=\operatorname{nrd}(I).
% \]
% We regard $I$ as a rank-four lattice with respect to the inner product
% \[
% \langle x,y\rangle=\operatorname{trd}(x\overline y),
% \]
% so that
% \[
% \|x\|^2=\langle x,x\rangle=2\operatorname{nrd}(x).
% \]
% Let
% \[
% I^\#=\{y\in B_{p,\infty}:\operatorname{trd}(x\overline y)\in\mathbb Z
% \text{ for all }x\in I\}
% \]
% be the trace-dual lattice.  Let $\alpha_1^\#\in I^\#$ be a shortest
% nonzero vector with respect to the reduced norm, and let
% $(\alpha_1,\ldots,\alpha_4)$ be a Minkowski-reduced basis of $I$, so that
% \[
% \operatorname{nrd}(\alpha_4)=\rho_4(I)
% \]
% is the fourth successive minimum measured in reduced norm.

% The rank-four transference theorem gives
% \[
% 1\le \|\alpha_4\|\cdot \|\alpha_1^\#\| \le 4.
% \]
% Since $\|x\|^2=2\operatorname{nrd}(x)$, this is equivalent to
% \[
% \frac14
% \le
% \operatorname{nrd}(\alpha_4)\operatorname{nrd}(\alpha_1^\#)
% \le
% 4.
% \]
% Dividing by $N=\operatorname{nrd}(I)$, we obtain
% \[
% \frac{1}{4\,N\operatorname{nrd}(\alpha_1^\#)}
% \le
% \frac{\operatorname{nrd}(\alpha_4)}{\operatorname{nrd}(I)}
% \le
% \frac{4}{N\operatorname{nrd}(\alpha_1^\#)}.
% \]

% Our experiment gives
% \[
% \log_2\!\left(N\operatorname{nrd}(\alpha_1^\#)\right)
% =
% -\frac12\log_2 p+O(1),
% \]
% or equivalently
% \[
% N\operatorname{nrd}(\alpha_1^\#)=\Theta(p^{-1/2}).
% \]
% Therefore, by transference,
% \[
% \frac{\operatorname{nrd}(\alpha_4)}{\operatorname{nrd}(I)}
% =
% \Theta(\sqrt p).
% \]
% More explicitly, if
% \[
% N\operatorname{nrd}(\alpha_1^\#)\ge p^{-1/2},
% \]
% then
% \[
% \frac{\operatorname{nrd}(\alpha_4)}{\operatorname{nrd}(I)}
% \le 4\sqrt p.
% \]
% Thus the observed slope $-1/2$ for
% \[
% \log_2\!\left(N\operatorname{nrd}(\alpha_1^\#)\right)
% \]
% is consistent with the heuristic expectation that the normalized fourth
% successive reduced norm of $I$ is of order $\sqrt p$.





% \subsection{Tighter analysis of quaternion vector sampling}

% We next refine the size analysis of the quaternion vector sampled by the algorithm \textsf{RandomEquivalentQuaternion} (Algorithm~\ref{alg:RandomEquivalentQuaternion}).
% In the previous analysis, the relevant sampling bound was estimated rather conservatively:
% the bottlenecks from HNF and \textsf{RandomIdealGivenNorm} increased the size bound strictly larger than $O(p^6)$.

% \begin{algorithm}
% \caption{RandomEquivalentQuaternion($L$)}
% \label{alg:RandomEquivalentQuaternion}
% \textbf{Input:} A lattice $L$.\\
% \textbf{Output:} A uniformly sampled $\mathbf{b}\in L$ such that
% $\operatorname{nrd}(\mathbf{b}) < D_{\mathrm{rsp}}\cdot D_{\mathrm{mix}}^{2}\cdot 2^{\,f}$.
% \begin{algorithmic}[1]
% \State Let $(b_{0},b_{1},b_{2},b_{3})$ be a basis of $L$
% \State Compute the Gram matrix $\mathbf{G}$ of $(b_{0},b_{1},b_{2},b_{3})$
% \State $\mathbf{b} \gets \textbf{LatticeSampling}\!\big((b_{0},b_{1},b_{2},b_{3}),\, \mathbf{G},\,
%       D_{\mathrm{rsp}}\cdot D_{\mathrm{mix}}^{2}\cdot 2^{\,f+1}\big)$
% \State \Return $\mathbf{b}$
% \end{algorithmic}
% \end{algorithm}

% The following lemma records the resulting tighter integer-size bound.
% The improvement is purely quantitative: the algorithm is unchanged, but the bound used in the complexity analysis is reduced from $O(p^6)$ to $O(p^{5.5})$.

% \begin{lemma}
% During the execution of Algorithm~\ref{alg:RandomEquivalentQuaternion}, the maximum size of integers is less than $16p^{5.5}$.
% \end{lemma}

% \begin{proof}
% The maximum size occurs in line~3.
% By~\cite[Lemma~25]{KLKL25}, the maximum size of integers is at most
% \[ \Drsp\cdot\Dmix^2\cdot2^{f+1} < 16p^{5.5}, \]
% since $\Drsp < 2\sqrt{p},\; \Dmix < 2p^2,\; 2^f < p$.
% \end{proof}